% Part: propositional-logic
% Chapter: syntax-and-semantics
% Section: formulas

\documentclass[../../../include/open-logic-section]{subfiles}

\begin{document}

\olfileid[fr]{pl}{syn}{fml}

\olsection{\usetoken{P}{formula} propositionnelles}

Les !!{formula}s de la logique propositionnelle sont construites à
partir de \emph{!!{propositional variable}s}%
\iftag{prvFalse}{\iftag{prvTrue}{,}{ et} de la constante
propositionnelle~$\lfalse$}{}%
\iftag{prvTrue}{ et de la constante
propositionnelle~$\ltrue$}{} au moyen de \emph{connecteurs logiques}.

\begin{enumerate}
\item Un ensemble~$\PVar$ !!{denumerable} de
  !!{propositional variable}s $\Obj p_0$, $\Obj p_1$, \dots
\tagitem{prvFalse}{La constante propositionnelle de
  !!{falsity}~$\lfalse$.}{}
\tagitem{prvTrue}{La constante propositionnelle de
  !!{truth}~$\ltrue$.}{}
\item Les connecteurs logiques :
  \startycommalist
  \iftag{prvNot}{\ycomma $\lnot$ (négation)}{}%
  \iftag{prvAnd}{\ycomma $\land$ (conjonction)}{}%
  \iftag{prvOr}{\ycomma $\lor$ (disjonction)}{}%
  \iftag{prvIf}{\ycomma $\lif$ (!!{conditional})}{}%
  \iftag{prvIff}{\ycomma $\liff$ (!!{biconditional})}{}%
\item Les signes de ponctuation : (, ) et la virgule.
\end{enumerate}

On note $\Lang L_0$ ce langage de la logique propositionnelle.

\iftag{defNot,defOr,defAnd,defIf,defIff,defTrue,defFalse,defEx,defAll}{%
Outre les connecteurs primitifs ci-dessus, on utilise les symboles
\emph{définis} suivants :
\startycommalist
  \iftag{defNot}{\ycomma $\lnot$ (négation)}{}%
  \iftag{defOr}{\ycomma $\lor$ (disjonction)}{}%
  \iftag{defAnd}{\ycomma $\land$ (conjonction)}{}%
  \iftag{defIf}{\ycomma $\lif$ (!!{conditional})}{}%
  \iftag{defIff}{\ycomma $\liff$ (!!{biconditional})}{}%
  \iftag{defFalse}{\ycomma $\lfalse$ (!!{falsity})}{}%
  \iftag{defTrue}{\ycomma $\ltrue$ (!!{truth})}{}.}{}

\begin{tagblock}{defNot,defOr,defAnd,defIf,defIff,defTrue,defFalse,defEx,defAll}
\begin{explain}
Un symbole défini ne fait pas officiellement partie du langage :
il est introduit comme abréviation informelle. Il permet d'abréger
des formules qui deviendraient assez longues si l'on n'utilisait que
les symboles primitifs. Cet avantage est évident, mais il y en a un
plus important : les démonstrations deviennent plus courtes. Tout
symbole primitif doit en effet y être traité séparément. Plus il y a
de symboles primitifs, plus les démonstrations sont longues.
\end{explain}
\end{tagblock}

% Alternate symbols

\begin{intro}
Vous connaissez peut-être d'autres termes ou symboles que ceux
employés ci-dessus. Les ouvrages de logique et les enseignants
utilisent couramment $\sim$, $\neg$ ou~!\ pour la « négation »,
et $\wedge$, $\cdot$ ou $\&$ pour la « conjonction ».
Le « conditionnel », ou l'« implication », se note souvent
$\rightarrow$, $\Rightarrow$ ou $\supset$.
\iftag{prvIff,defIff}{Le « biconditionnel », la « double implication »
ou l'« équivalence (matérielle) » se notent $\leftrightarrow$,
$\Leftrightarrow$ ou $\equiv$.}{}
\iftag{prvFalse,defFalse}{Le symbole $\lfalse$ est diversement appelé
« faux », « falsum », « absurdité » ou « bottom » (« bas »).}{}
\iftag{prvTrue,defTrue}{Le symbole $\ltrue$ est diversement appelé
« vrai », « verum » ou « top » (« haut »).}{}
\end{intro}

\begin{defn}[Formule]
\ollabel{defn:formulas}
L'ensemble~$\Frm[L_0]$ des \emph{!!{formula}s} de la logique
propositionnelle est défini inductivement comme suit :
\begin{enumerate}
\tagitem{prvFalse}{$\lfalse$ est une !!{formula} atomique.}{}

\tagitem{prvTrue}{$\ltrue$ est une !!{formula} atomique.}{}

\item Toute !!{propositional variable}~$\Obj p_i$ est une
  !!{formula} atomique.

\tagitem{prvNot}{Si $!A$ est une !!{formula}, alors $\lnot !A$
  est une !!{formula}.}{}

\tagitem{prvAnd}{Si $!A$ et $!B$ sont des !!{formula}s,
  alors $(!A \land !B)$ est une !!{formula}.}{}

\tagitem{prvOr}{Si $!A$ et $!B$ sont des !!{formula}s,
  alors $(!A \lor !B)$ est une !!{formula}.}{}

\tagitem{prvIf}{Si $!A$ et $!B$ sont des !!{formula}s,
  alors $(!A \lif !B)$ est une !!{formula}.}{}

\tagitem{prvIff}{Si $!A$ et $!B$ sont des !!{formula}s,
  alors $(!A \liff !B)$ est une !!{formula}.}{}

\tagitem{limitClause}{Rien d'autre n'est une !!{formula}.}{}
\end{enumerate}
\end{defn}

\begin{explain}
La définition des !!{formula}s est une \emph{définition inductive}.
On construit, en substance, leur ensemble en une infinité d'étapes.
À l'étape initiale, on déclare que toutes les formules atomiques sont
des formules. Cela correspond aux premiers cas de la définition,
c'est-à-dire aux cas
\iftag{prvTrue}{$\ltrue$, }{}%
\iftag{prvFalse}{$\lfalse$, }{}%
$\Obj p_i$. Une « !!{formula} atomique » est donc une !!{formula}
de l'une de ces formes.

Les autres cas donnent des règles pour construire de nouvelles
!!{formula}s à partir de celles déjà construites. À la deuxième étape,
on applique ces règles aux !!{formula}s atomiques. À la troisième,
on construit de nouvelles formules à partir des formules atomiques
et de celles obtenues à la deuxième étape, et ainsi de suite.
Est une !!{formula} tout ce qui finit par être construit à l'une
de ces étapes, et rien d'autre.
\end{explain}

Lorsqu'on écrit une formule $(!B \ast !C)$ construite à partir de
$!B$, $!C$ au moyen d'un connecteur binaire~$\ast$, on omet souvent
la paire de parenthèses extérieure et l'on écrit simplement~$!B \ast !C$.

\begin{tagblock}{defNot,defOr,defAnd,defIf,defIff,defTrue,defFalse,defEx,defAll}
\begin{defn}
Les formules construites au moyen des opérateurs définis s'entendent
comme suit :

\begin{tagenumerate}{defTrue,defFalse,defNot,defOr,defAnd,defIf,defIff,defEx,defAll}

\tagitem{defTrue}{$\ltrue$ abrège
  \iftag{prvFalse}{$\lnot\lfalse$}{$(!A \lor \lnot !A)$, où
    l'on fixe une !!{formula} atomique~$!A$}.}{}

\tagitem{defFalse}{$\lfalse$ abrège
  \iftag{prvTrue}{$\lnot\ltrue$}{$(!A \land \lnot !A)$, où
    l'on fixe une !!{formula} atomique~$!A$}.}{}

\tagitem{defNot}{$\lnot !A$ abrège $!A \lif \lfalse$.}{}

\tagitem{defOr}{$!A \lor !B$ abrège
  \iftag{prvAnd}{$\lnot(\lnot !A \land \lnot !B)$}{$\lnot !A \lif
    !B$}.}{}

\tagitem{defAnd}{$!A \land !B$ abrège
  \iftag{prvOr}{$\lnot(\lnot !A \lor \lnot !B)$}{$\lnot (!A \lif
    \lnot !B)$}.}{}

\tagitem{defIf}{$!A \lif !B$ abrège
  \iftag{prvOr}{$(\lnot !A \lor !B)$}{$\lnot (!A \land \lnot !B)$}.}{}

\tagitem{defIff}{$!A \liff !B$ abrège $(!A \lif !B) \land (!B
  \lif !A)$.}{}
\end{tagenumerate}
\end{defn}
\begin{editorial}
Dans la branche où la disjonction est primitive, la source omet la
parenthèse ouvrante de l'abréviation de l'implication. Elle est rétablie ici.
\end{editorial}
\end{tagblock}

\begin{defn}[Identité syntaxique]
Le symbole $\ident$ exprime l'identité syntaxique de deux chaînes de
symboles : $!A \ident !B$ si et seulement si $!A$ et $!B$ ont la même
longueur et comportent le même symbole à chaque position.
\end{defn}

De part et d'autre du symbole $\ident$, on peut écrire des chaînes
obtenues par concaténation. Par exemple, $!A \ident (!B \lor !C)$
signifie que la chaîne~$!A$ est celle obtenue en concaténant, dans cet
ordre, une parenthèse ouvrante, la chaîne $!B$, le symbole $\lor$,
la chaîne~$!C$ et une parenthèse fermante. On sait alors que le premier
symbole de $!A$ est une parenthèse ouvrante, que $!A$ contient $!B$
comme sous-chaîne à partir du deuxième symbole, que cette sous-chaîne
est suivie de $\lor$, etc.

\end{document}
